\documentclass[a4paper,12pt]{article}

\usepackage{iftex}
\ifPDFTeX
  \usepackage[T2A]{fontenc}
  \usepackage[utf8]{inputenc}
  \usepackage[russian]{babel}
  \usepackage{helvet}
  \renewcommand{\familydefault}{\sfdefault}
\else
  \usepackage{fontspec}
  \usepackage{polyglossia}
  \setdefaultlanguage{russian}
  \setmainfont{DejaVu Sans}
  \setsansfont{DejaVu Sans}
  \renewcommand{\familydefault}{\sfdefault}
\fi

\usepackage{amsmath,amssymb,mathtools}
\usepackage{icomma}
\usepackage{geometry}
\usepackage{microtype}
\usepackage{tikz}
\usepackage{tabularx,booktabs,longtable,array}
\usepackage{enumitem}
\usepackage{hyperref}
\usepackage{parskip}
\usepackage{fancyhdr}
\usepackage{setspace}
\usepackage{xcolor}

\geometry{left=18mm,right=18mm,top=18mm,bottom=20mm}
\onehalfspacing
\setlength{\parindent}{0pt}
\setlength{\parskip}{6pt}
\setlength{\emergencystretch}{2em}
\setlist[itemize]{leftmargin=7mm,itemsep=3pt,topsep=4pt}
\setlist[enumerate]{leftmargin=8mm,itemsep=4pt,topsep=4pt}
\renewcommand{\arraystretch}{1.35}

\definecolor{accent}{HTML}{7A3B45}
\definecolor{darkaccent}{HTML}{4D2730}
\definecolor{textgray}{HTML}{333333}
\definecolor{softgray}{HTML}{6E6E6E}
\definecolor{paperline}{HTML}{D8C8BC}

\hypersetup{colorlinks=true,linkcolor=darkaccent,urlcolor=accent,pdftitle={Финансовая математика: вклады и сложные проценты},pdfauthor={Лёвин Артём Александрович}}
\pagestyle{fancy}
\fancyhf{}
\lhead{\textcolor{textgray}{Финансовая математика}}
\rhead{\textcolor{textgray}{Вклады и сложные проценты}}
\cfoot{\textcolor{softgray}{\thepage}}
\setlength{\headheight}{15pt}

\newcommand{\term}[1]{\textcolor{darkaccent}{\textbf{#1}}}
\newcommand{\formula}[1]{%
  \begin{center}
    \begin{minipage}{0.92\linewidth}
      \centering\large\color{darkaccent}#1
    \end{minipage}
  \end{center}}
\newcommand{\checkitem}{\item[$\square$]}

\begin{document}

\begin{titlepage}
  \thispagestyle{empty}
  \centering
  \vspace*{26mm}
  {\Large\textcolor{accent}{Чек-лист}\par}
  \vspace{7mm}
  {\Huge\bfseries\textcolor{darkaccent}{Финансовая математика}\par}
  \vspace{3mm}
  {\LARGE Вклады, сложные проценты и дополнительные операции\par}
  \vfill
  {\large Лёвин Артём Александрович\par}
  {\large эксклюзивно для Николь Саркисянц\par}
  \vspace{8mm}
  {\large \today\par}
  \vfill
  \begin{tikzpicture}[x=1cm,y=1cm]
    \draw[accent,line width=1.2pt]
      (-7,0) .. controls (-4.7,1.25) and (-2.3,-1.15) .. (0,0)
             .. controls (2.3,1.15) and (4.7,-1.25) .. (7,0);
    \draw[paperline,line width=.6pt]
      (-7,-.35) .. controls (-3.5,.45) and (3.5,-.45) .. (7,-.35);
  \end{tikzpicture}
  \vspace*{8mm}
\end{titlepage}

\tableofcontents
\clearpage

\section{Язык финансовой модели}

Задача о вкладе описывает последовательность периодов. В каждом периоде важно знать порядок событий: какая сумма была на счёте, сколько начислил банк и какая операция выполнена после начисления.

\begin{table}[ht]
\centering
\caption{Основные обозначения}
\begin{tabularx}{\linewidth}{@{}>{\bfseries}lX@{}}
\toprule
Обозначение & Смысл \\
\midrule
$S_0$ & первоначальная сумма вклада; \\
$P\%$ & процентная ставка за один период; \\
$p=\dfrac{P}{100}$ & ставка в виде десятичной дроби; \\
$r=1+p$ & коэффициент роста за один период; \\
$A_k$ & дополнительная операция после начисления процентов в периоде $k$: взнос положителен, снятие отрицательно; \\
$S_k$ & сумма на счёте после всех действий периода $k$. \\
\bottomrule
\end{tabularx}
\end{table}

\term{Ключевая рекуррентная формула:}
\formula{$S_k=rS_{k-1}+A_k$.}

Сначала начисляются проценты, затем учитывается $A_k$, если условие задаёт именно такой порядок. Формулировку задачи следует читать буквально: изменение порядка действий меняет результат.

\subsection{Таблица одного периода}

\begin{table}[ht]
\centering
\caption{Универсальная строка расчёта}
\begin{tabularx}{\linewidth}{@{}cXXXX@{}}
\toprule
Период & Сумма в начале & Начисленные проценты & Операция & Сумма в конце \\
\midrule
$k$ & $S_{k-1}$ & $pS_{k-1}$ & $A_k$ & $S_{k-1}+pS_{k-1}+A_k=rS_{k-1}+A_k$ \\
\bottomrule
\end{tabularx}
\end{table}

\section{Сложные проценты}

Пусть $A_k=0$ в каждом периоде. Тогда следующая сумма получается умножением предыдущей на $r$:
\[
S_1=S_0r,\qquad S_2=S_0r^2,\qquad S_3=S_0r^3.
\]
Закономерность приводит к формуле сложных процентов:
\formula{$\boxed{S_n=S_0r^n=S_0\left(1+\dfrac{P}{100}\right)^n}$.}

\term{Почему формула работает.} За каждый период вся накопленная сумма умножается на один и тот же коэффициент $r$. После $n$ периодов этот множитель применяется $n$ раз.

\subsection{Пример 1. Два периода}

На вклад положили 120 рублей под 15\% за период. Дополнительных операций нет.
\[
p=0{,}15,\qquad r=1{,}15.
\]

\begin{table}[ht]
\centering
\caption{Расчёт вклада}
\begin{tabularx}{\linewidth}{@{}cXXXX@{}}
\toprule
Период & Начало & Проценты & Операция & Конец \\
\midrule
1 & $120$ & $18$ & $0$ & $138$ \\
2 & $138$ & $20{,}7$ & $0$ & $158{,}7$ \\
\bottomrule
\end{tabularx}
\end{table}

Проверка по формуле: $120\cdot1{,}15^2=158{,}7$.

\section{Удвоение вклада и округление}

Если вклад должен увеличиться в $q$ раз, то
\[
qS_0=S_0r^n.
\]
При $S_0>0$ можно сократить обе части на $S_0$:
\formula{$q=r^n$.}

\subsection{Найти срок}

При известной ставке точное значение срока выражается через логарифм:
\formula{$n=\dfrac{\ln q}{\ln r}$.}

Для удвоения при ставке 10\%:
\[
n=\frac{\ln2}{\ln1{,}1}\approx7{,}27.
\]
Требуется количество \emph{полных} лет, к концу которых вклад уже удвоится. Семь лет недостаточно, поэтому ответ равен 8 годам. На экзамене тот же вывод можно получить последовательной проверкой целых значений $n$.

\subsection{Найти ставку}

Если вклад удваивается за $n$ периодов, то
\[
2=r^n,\qquad r=\sqrt[n]{2},\qquad P=100\left(\sqrt[n]{2}-1\right).
\]
За пять лет:
\[
P=100\left(\sqrt[5]{2}-1\right)\approx14{,}87\%.
\]
Если ставка должна быть целым числом процентов и гарантировать удвоение, выбирается 15\%.

\term{Правило округления.} Направление определяет смысл условия.

Минимальный целый срок округляют вверх. Целую ставку выбирают так, чтобы требуемый результат уже был достигнут.

\section{Взносы и снятия}

При операции после начисления процентов используется формула
\[
S_k=rS_{k-1}+A_k.
\]
Положительное $A_k$ означает пополнение, отрицательное --- снятие.

\subsection{Пример 2. Постоянное пополнение}

На счёте 100 рублей, ставка 10\%, после каждого начисления добавляют 20 рублей.

\begin{table}[ht]
\centering
\caption{Три периода с пополнением}
\begin{tabularx}{\linewidth}{@{}cXXXX@{}}
\toprule
Период & Начало & Проценты & Взнос & Конец \\
\midrule
1 & $100$ & $10$ & $+20$ & $130$ \\
2 & $130$ & $13$ & $+20$ & $163$ \\
3 & $163$ & $16{,}3$ & $+20$ & $199{,}3$ \\
\bottomrule
\end{tabularx}
\end{table}

\subsection{Общая формула при постоянном взносе}

Если после каждого периода вносится одна и та же сумма $A$, то
\[
S_1=S_0r+A,
\]
\[
S_2=S_0r^2+Ar+A,
\]
\[
S_3=S_0r^3+Ar^2+Ar+A.
\]
Следовательно,
\formula{$S_n=S_0r^n+A\left(r^{n-1}+r^{n-2}+\dots+r+1\right)$.}

При $r\ne1$ сумму геометрической прогрессии можно записать компактно:
\formula{$S_n=S_0r^n+A\dfrac{r^n-1}{r-1}$.}

Эту формулу полезно понимать через таблицу. В развёрнутом решении экзаменационной задачи таблица или рекуррентная запись показывает происхождение каждого слагаемого.

\section{Сравнение двух сценариев}

Вклад открыт под 10\% годовых. После первого начисления вкладчик снял 2000 рублей, после второго начисления внёс 2000 рублей. Нужно определить, насколько фактическая сумма через три года меньше суммы без операций.

Пусть $r=1{,}1$ и первоначальный вклад равен $S_0$.

\term{План без операций:}
\[
S_{\text{план}}=S_0r^3.
\]

\term{Фактический сценарий:}
\[
S_1=S_0r-2000,
\]
\[
S_2=(S_0r-2000)r+2000,
\]
\[
S_3=\bigl((S_0r-2000)r+2000\bigr)r.
\]

Разность не зависит от $S_0$:
\begin{align*}
S_{\text{план}}-S_3
&=S_0r^3-\left(S_0r^3-2000r^2+2000r\right)\\
&=2000r^2-2000r\\
&=2000\cdot1{,}1\cdot0{,}1=220.
\end{align*}

\term{Ответ:} фактическая сумма меньше плановой на 220 рублей.

\section{Универсальный алгоритм}

\begin{enumerate}
  \item Выделите период начисления: год, месяц или другой интервал.
  \item Запишите ставку в виде десятичной дроби $p=P/100$.
  \item Введите коэффициент роста $r=1+p$.
  \item Уточните порядок событий внутри каждого периода.
  \item Составьте таблицу: сумма в начале, проценты, операция, сумма в конце.
  \item Перенесите конечную сумму периода в начало следующего.
  \item Получите формулу или уравнение и упростите его вынесением общих множителей.
  \item Решите уравнение, соблюдая условия $S_0>0$, $r>0$ и целочисленность срока, если она требуется.
  \item Проверьте направление округления по смыслу вопроса.
  \item Сравните результат с исходными данными и укажите денежные единицы.
\end{enumerate}

\section{Типичные ошибки}

\begin{itemize}
  \item Начислять каждый раз один и тот же процент от первоначальной суммы. Проценты считаются от текущего остатка.
  \item Перепутать $P$ и $p$: например, подставить 10 вместо $0{,}1$.
  \item Поменять местами начисление процентов и дополнительную операцию.
  \item Потерять знак операции: снятие записывается отрицательным числом.
  \item Раскрыть скобки раньше времени и допустить арифметическую ошибку.
  \item Округлить срок вниз, хотя условие требует уже достигнутого результата.
  \item Сравнивать сценарии с разным числом периодов или разными начальными суммами.
\end{itemize}

\section{Итоговый чек-лист}

\begin{itemize}
  \checkitem Я расшифровал(а) все введённые обозначения.
  \checkitem Я перевёл(а) процентную ставку в десятичную дробь.
  \checkitem Я записал(а) коэффициент роста $r=1+p$.
  \checkitem Я установил(а) порядок начислений, взносов и снятий.
  \checkitem Я переношу конечную сумму периода в начало следующего.
  \checkitem Я использую рекуррентную формулу $S_k=rS_{k-1}+A_k$.
  \checkitem Я могу вывести формулу $S_n=S_0r^n$.
  \checkitem Я понимаю происхождение каждого слагаемого при постоянных взносах.
  \checkitem Я выбираю направление округления по смыслу задачи.
  \checkitem Я проверяю ответ сравнением сценариев и здравым смыслом.
\end{itemize}

\clearpage
\section{Тренировочный блок}
\begin{sloppypar}

\textbf{Средний уровень}

\begin{enumerate}
  \item Вклад 50\,000 рублей открыт под 8\% годовых. Найдите сумму через два года, если операций со счётом нет.
  \item На счёт внесли 120\,000 рублей под 10\% годовых. После каждого начисления процентов добавляют 10\,000 рублей. Найдите сумму через два года.
  \item Вклад 90\,000 рублей открыт под 12\% годовых. После первого начисления сняли 15\,000 рублей. Найдите сумму после второго начисления, если других операций нет.
\end{enumerate}

\textbf{Уровень выше среднего}

\begin{enumerate}[resume]
  \item Найдите минимальное число полных лет, за которое вклад увеличится не менее чем в три раза при ставке 15\% годовых.
  \item Найдите минимальную целую процентную ставку, при которой вклад за четыре года увеличится не менее чем вдвое.
  \item На счёт внесли 100\,000 рублей под 6\% годовых. После каждого начисления добавляют 5000 рублей. Найдите сумму через три года.
  \item Вклад открыт под 8\% годовых. После первого начисления сняли 3000 рублей, после второго внесли 3000 рублей. На сколько фактическая сумма через три года меньше плановой?
  \item После первого начисления 10\% к вкладу добавили 20\,000 рублей, после второго начисления сняли 10\,000 рублей. В итоге на счёте оказалось 133\,000 рублей. Найдите первоначальную сумму.
  \item При ставке 10\% после каждого из трёх годовых начислений вносили одинаковую сумму $A$. Через три года итог превысил сумму без пополнений на 33\,100 рублей. Найдите $A$.
\end{enumerate}

\textbf{Сложный уровень}

\begin{enumerate}[resume]
  \item Вклад открыт под целое число процентов годовых. После первого начисления сняли 20\,000 рублей, после второго внесли 10\,000 рублей, после третьего операций не было. Через три года фактическая сумма оказалась на 16\,800 рублей меньше плановой. Найдите процентную ставку.
\end{enumerate}
\end{sloppypar}

\clearpage
\section{Ответы}

\begin{table}[ht]
\centering
\caption{Ответы к тренировочному блоку}
\begin{tabularx}{\linewidth}{@{}cX@{}}
\toprule
№ & Ответ \\
\midrule
1 & 58\,320 рублей. \\
2 & 166\,200 рублей. \\
3 & 96\,096 рублей. \\
4 & 8 лет. \\
5 & 19\%. \\
6 & 135\,019{,}60 рубля. \\
7 & 259{,}20 рубля. \\
8 & 100\,000 рублей. \\
9 & 10\,000 рублей. \\
10 & 20\%. \\
\bottomrule
\end{tabularx}
\end{table}

\vfill
\begin{center}
\textcolor{softgray}{Код самопроверки: 1426}
\end{center}

\end{document}
